\begin{tikzpicture}[
  font=\scriptsize,
  panel/.style={draw, rounded corners, fill=white, minimum width=2.8cm,
    minimum height=3.1cm},
  judul/.style={font=\scriptsize\bfseries, align=center},
  infra/.style={draw, rounded corners, fill=blue!8, align=center,
    text width=2.2cm, minimum height=0.75cm, font=\tiny},
  infrakecil/.style={draw, rounded corners, fill=blue!8, align=center,
    text width=0.7cm, minimum height=0.75cm, font=\tiny},
  org/.style={draw, fill=praditagreen!15, align=center, minimum width=0.55cm,
    minimum height=0.4cm, font=\tiny},
  orglebar/.style={draw, fill=praditagreen!15, align=center,
    text width=2.2cm, minimum height=0.4cm, font=\tiny},
  catatan/.style={font=\tiny\mdseries, text=gray!80!black, align=center,
    text width=2.6cm},
  garis/.style={thick, gray!70!black},
  panah/.style={-{Stealth[length=2mm]}, thick}
]
  \node[font=\scriptsize\bfseries] at (4.5,0.45)
    {Empat model penerapan menurut NIST SP 800-145};

  % Public cloud
  \node[panel] at (0,-1.45) {};
  \node[judul] at (0,-0.2) {\textit{Public cloud}};
  \node[infra] (pubinfra) at (0,-0.95) {Infrastruktur milik penyedia};
  \node[org] (puba) at (-0.75,-1.95) {A};
  \node[org] (pubb) at (0,-1.95) {B};
  \node[org] (pubc) at (0.75,-1.95) {C};
  \foreach \o in {puba, pubb, pubc} {\draw[garis] (\o.north) -- (\o.north |- pubinfra.south);}
  \node[catatan] at (0,-2.6) {terbuka untuk umum, dipakai banyak organisasi};

  % Private cloud
  \node[panel] at (3,-1.45) {};
  \node[judul] at (3,-0.2) {\textit{Private cloud}};
  \node[infra] (priinfra) at (3,-0.95) {Infrastruktur khusus satu organisasi};
  \node[orglebar] (prio) at (3,-1.95) {Satu organisasi};
  \draw[garis] (prio.north) -- (priinfra.south);
  \node[catatan] at (3,-2.6) {di lokasi sendiri atau dikelola pihak ketiga};

  % Community cloud
  \node[panel] at (6,-1.45) {};
  \node[judul] at (6,-0.2) {\textit{Community cloud}};
  \node[infra] (kominfra) at (6,-0.95) {Infrastruktur bersama satu komunitas};
  \node[org] (koma) at (5.25,-1.95) {A};
  \node[org] (komb) at (6,-1.95) {B};
  \node[org] (komc) at (6.75,-1.95) {C};
  \foreach \o in {koma, komb, komc} {\draw[garis] (\o.north) -- (\o.north |- kominfra.south);}
  \node[catatan] at (6,-2.6) {kepentingan sama, misalnya syarat kepatuhan};

  % Hybrid cloud
  \node[panel] at (9,-1.45) {};
  \node[judul] at (9,-0.2) {\textit{Hybrid cloud}};
  \node[infrakecil] (hibpri) at (8.35,-0.95) {Privat};
  \node[infrakecil] (hibpub) at (9.65,-0.95) {Publik};
  \draw[{Stealth[length=1.5mm]}-{Stealth[length=1.5mm]}, thick]
    (hibpri.east) -- (hibpub.west);
  \node[orglebar] (hibo) at (9,-1.95) {Satu organisasi};
  \draw[garis] (hibo.north -| hibpri.south) -- (hibpri.south);
  \draw[garis] (hibo.north -| hibpub.south) -- (hibpub.south);
  \node[catatan] at (9,-2.6) {dua infrastruktur terpisah, data dan aplikasi dapat berpindah};

  % Multi-cloud
  \node[panel, minimum width=5.6cm, minimum height=2.6cm, fill=orange!12]
    at (4.5,-4.75) {};
  \node[judul] at (4.5,-3.75)
    {\textit{Multi-cloud} {\tiny\mdseries (praktik industri, bukan model NIST)}};
  \node[infra, text width=1.9cm] (mca) at (3.3,-4.45) {Penyedia publik A};
  \node[infra, text width=1.9cm] (mcb) at (5.7,-4.45) {Penyedia publik B};
  \node[orglebar, text width=3.0cm] (mco) at (4.5,-5.45) {Satu organisasi};
  \draw[garis] (mco.north -| mca.south) -- (mca.south);
  \draw[garis] (mco.north -| mcb.south) -- (mcb.south);
\end{tikzpicture}
